% !TEX root = ../../thesis.tex
\section{Settings and Preliminaries}\label{ch:preliminaries}

This section establishes the foundational framework for the free boundary min-max theory in complete Riemannian manifolds with boundary.

Let $(M^{n+1}, g)$ be a smooth, complete Riemannian manifold with non-empty smooth boundary $\partial M$. We assume $2 \leq n \leq 6$ throughout. We write $\nu_{\partial M}$ for the inward unit normal to $\partial M$. For details on Fermi coordinates near $\partial M$, see \cite[Section~2]{LZ17}.

We denote by $\mathcal{H}^k$ the $k$-dimensional Hausdorff measure and by $\mathbf{M}(T)$ the mass of a current $T$. We assume $M$ is isometrically embedded in some $\mathbb{R}^L$. For vector fields, let $\mathfrak{X}_c(M)$ denote the space of \emph{compactly supported} smooth vector fields on $M$ that are tangent to $M$:
\begin{multline*}
\mathfrak{X}_c(M) := \bigl\{X \in \mathfrak{X}(\mathbb{R}^L) : X \text{ is compactly supported},\\
X(p) \in T_p M \text{ for all } p \in M\bigr\}.
\end{multline*}
The subspace of vector fields additionally tangent to the boundary is
\[
\mathfrak{X}_{\mathrm{tan}}(M) := \{X \in \mathfrak{X}_c(M) : X(p) \in T_p(\partial M) \text{ for all } p \in \partial M\}.
\]

\begin{remark}[Role of the isometric embedding]\label{rem:isometric-embedding}
The assumption that $M$ is isometrically embedded in $\mathbb{R}^L$ (which exists by Nash's embedding theorem~\cite{Nash56}) is needed because integer rectifiable currents are defined as functionals on differential forms in an ambient Euclidean space (see Simon~\cite[\S26--27]{Sim84}, Li--Zhou~\cite[\S3]{LZ17}).
\end{remark}

For a relatively open subset $U \subset M$, we write $\mathrm{int}_M(U)$ for the relative interior of $U$ in $M$. We use $A \Subset B$ to denote that $\overline{A}$ is compact and contained in $B$.


%%%%%%%%%%%%%%%%%%%%%%%%%%%%%%%%%%%%%%%%%%%%%%%%%%%%%%%%%%%%%%%%%%%%%%%%%%%%%%%
\subsection{Currents and Varifolds}\label{sec:currents-varifolds}
%%%%%%%%%%%%%%%%%%%%%%%%%%%%%%%%%%%%%%%%%%%%%%%%%%%%%%%%%%%%%%%%%%%%%%%%%%%%%%%

We work with currents and relative cycles following Li--Zhou~\cite{LZ17}. We write $\mathcal{R}_k(M;\mathbb{Z}_2)$ for the space of mod~$2$ flat chains supported in $M$, with mass norm $\mathbf{M}$ and flat topology.

\begin{remark}[Coefficient group]\label{rmk:Z2}
Throughout this thesis we work with $\mathbb{Z}_2$ coefficients. The reduced boundary $\partial^*\Omega$ of a Caccioppoli set carries no canonical orientation, only a distinction between interior and exterior, so $\mathbb{Z}_2$ is the natural coefficient group. Li--Zhou~\cite{LZ17} state their results for $\mathbb{Z}$ coefficients, but the key tools---the isoperimetric lemma, slicing theorem, and regularity theory---hold equally for $\mathbb{Z}_2$ (see \cite{MN17,Alm62}).
\end{remark}

\begin{definition}[Relative cycles and equivalence classes]\label{def:relative-cycle}
For subsets $A \supset B$ of $M$, the space of \emph{relative cycles} is
\[
\mathcal{Z}_n(A, B;\mathbb{Z}_2) := \{T \in \mathcal{R}_n(M;\mathbb{Z}_2) : \mathrm{spt}(T) \subset A, \ \mathrm{spt}(\partial T) \subset B\},
\]
equipped with the flat topology. The \emph{space of equivalence classes of relative cycles} is the quotient
\[
\mathbf{Z}_n(A, B;\mathbb{Z}_2) := \mathcal{Z}_n(A, B;\mathbb{Z}_2) / \mathcal{Z}_n(B, B;\mathbb{Z}_2).
\]
Two cycles $T_1, T_2 \in \mathcal{Z}_n(A, B;\mathbb{Z}_2)$ represent the same class $\tau \in \mathbf{Z}_n(A, B;\mathbb{Z}_2)$ if and only if $T_1 - T_2 \in \mathcal{Z}_n(B, B;\mathbb{Z}_2)$, i.e., $T_1 - T_2$ is supported in $B$.
\end{definition}

\begin{proposition}[Caccioppoli sets and mod~$2$ cycles]\label{prop:cacc-cycle}
Let $\Omega \subset M$ be a measurable subset. Then $\Omega$ is a Caccioppoli set (i.e.\ $\mathrm{Per}(\Omega) < \infty$) if and only if the integration current $\llrr{\Omega} \in \mathcal{R}_{n + 1}(M; \Z_2)$ satisfies $\mass{(\partial\llrr{\Omega})} < \infty$. In this case,
\begin{equation*}
    \partial\llrr{\Omega} \in \mathcal{Z}_n(M, \partial M;\Z_2)
\end{equation*}
is a mod~$2$ relative cycle with
\begin{equation*}
    \mathrm{Per}(\Omega) = \mass{(\partial\llrr{\Omega})}.
\end{equation*}
The current $\partial\llrr{\Omega}$ is supported on the reduced boundary: $\partial^*\Omega \subset \spt{\partial\llrr{\Omega}}\subset\overline{\partial^*\Omega}$, and the varifold $\abs{\partial^*\Omega}$ induced by the reduced boundary agrees with the varifold $\abs{\partial\llrr{\Omega}}$ induced by the mod~$2$ current. Moreover, $\Omega$ and $M\setminus \Omega$ induce the same mod~$2$ cycle: $\partial\llrr{\Omega} = \partial\llrr{M\setminus \Omega}$.
\end{proposition}
\begin{proof}
The set $\Omega$ defines an integer rectifiable $(n{+}1)$-current $\llrr{\Omega}$ by integration: $\llrr{\Omega}(\omega) = \int_\Omega \omega$ for any smooth $(n{+}1)$-form $\omega$. Reducing modulo~$2$, we regard $\llrr{\Omega} \in \mathcal{R}_{n+1}(M; \Z_2)$. By definition, $\Omega$ is a Caccioppoli set if and only if its distributional boundary $\mu_\Omega$ is a finite Radon measure, i.e.\ $\abs{\mu_\Omega}(M) < \infty$. Since both $\abs{\mu_\Omega}(M)$ and $\mass{(\partial\llrr{\Omega})}$ express the total variation of the distributional derivative of $\mathbf{1}_\Omega$, we have $\abs{\mu_\Omega}(M) = \mass{(\partial\llrr{\Omega})}$, and the equivalence $\mathrm{Per}(\Omega) < \infty \Leftrightarrow \mass{(\partial\llrr{\Omega})} < \infty$ follows.

The chain complex axiom gives $\partial(\partial\llrr{\Omega}) = \partial^2\llrr{\Omega} = 0$, so $\mathrm{spt}(\partial(\partial\llrr{\Omega})) = \emptyset \subset \partial M$ and $\partial\llrr{\Omega} \in \mathcal{Z}_n(M, \partial M; \Z_2)$.

By De~Giorgi's structure theorem (Maggi~\cite[Theorem~15.9]{Mag12}), the Gauss--Green measure satisfies $\abs{\mu_\Omega} = \mathcal{H}^n \llcorner \partial^*\Omega$, and the density $\Theta^n(\partial^*\Omega, x) = 1$ for every $x \in \partial^*\Omega$ (Maggi~\cite[Corollary~15.8]{Mag12}). Therefore
\[
\mass{(\partial\llrr{\Omega})} = \abs{\mu_\Omega}(M) = \mathcal{H}^n(\partial^*\Omega) = \mathrm{Per}(\Omega),
\]
and $\partial^*\Omega \subset \spt{\partial\llrr{\Omega}} \subset \overline{\partial^*\Omega}$. Note that the current $\partial\llrr{\Omega}$ is defined by first forming the integration current and then applying $\partial$, requiring no a priori knowledge of the reduced boundary. One could alternatively define an $n$-current $\llrr{\partial^*\Omega}$ by integrating directly over $\partial^*\Omega$; De~Giorgi's structure theorem identifies the two: $\partial\llrr{\Omega} = \llrr{\partial^*\Omega}$. Throughout this thesis we write $\llbracket\partial\Omega\rrbracket := \partial\llrr{\Omega}$, with the understanding that $\partial\Omega$ refers to the reduced boundary $\partial^*\Omega$ in the context of currents.

Since $\partial\llrr{\Omega}$ is carried by $\partial^*\Omega$ with multiplicity one, the associated varifold $\abs{\partial\llrr{\Omega}}$ acts on $f \in C_c(G_n(M))$ by
\[
\abs{\partial\llrr{\Omega}}(f) = \int_{\partial^*\Omega} f(x, T_x(\partial^*\Omega)) \, d\mathcal{H}^n(x) = \abs{\partial^*\Omega}(f),
\]
where $T_x(\partial^*\Omega)$ is the approximate tangent plane, which exists $\mathcal{H}^n$-a.e.\ by rectifiability.

Finally, in $\Z_2$ coefficients we have $\llrr{\Omega} + \llrr{M \setminus \Omega} = \llrr{M}$. Since $\partial\llrr{M}$ is supported on $\partial M$, it vanishes as a relative cycle in $\mathcal{Z}_{n}(M, \partial M; \Z_2)$, so $\partial\llrr{\Omega} + \partial\llrr{M \setminus \Omega} = 0$, which in $\Z_2$ gives $\partial\llrr{\Omega} = \partial\llrr{M \setminus \Omega}$.
\end{proof}

The min-max machinery of Li--Zhou~\cite{LZ17} operates on equivalence classes of relative cycles rather than on individual currents. We therefore need a well-defined notion of mass on these classes.

\begin{definition}[Canonical representative]\label{def:canonical-rep}
For $\tau \in \mathbf{Z}_n(A, B;\mathbb{Z}_2)$, there exists a unique representative $T \in \tau$ satisfying $T \llcorner B = 0$. This $T$ is called the \emph{canonical representative} of $\tau$. When no confusion arises, we identify $\tau$ with its canonical representative.
\end{definition}

\begin{definition}[Mass of an equivalence class]\label{def:mass-equivalence-class}
For $\tau \in \mathbf{Z}_n(A, B;\mathbb{Z}_2)$, the \emph{mass} of $\tau$ is defined as $\mathbf{M}(\tau) := \mathbf{M}(T)$, where $T$ is the canonical representative of $\tau$. By \cite[Lemma~3.3]{LZ17}, $\mathbf{M}(\tau) = \inf_{S \in \tau} \mathbf{M}(S)$, so the canonical representative achieves the infimum of mass over all representatives.
\end{definition}

Given a Caccioppoli set $\Omega$, we refer to the $(n{+}1)$-current $\llrr{\Omega}$ as the \emph{filling} of the relative cycle $\partial\llrr{\Omega}$. For $\tau \in \mathbf{Z}_n(M, \partial M;\mathbb{Z}_2)$ with canonical representative $T$, the mass restricted to an open set $V$ is $\|\tau\|(V) := \mathbf{M}(T \llcorner V)$.

\begin{definition}[Interior boundary and free boundary part]\label{def:boundary-decomposition}
For a measurable set $\Omega \subset M$ with $\mass{(\partial\llrr{\Omega})} < \infty$, the boundary current decomposes as
\[
\partial \llrr{\Omega} = \partial_I \Omega + \partial_F \Omega,
\]
where:
\begin{itemize}
    \item $\partial_I \Omega := \partial \llrr{\Omega} \llcorner \mathrm{int}(M)$ is the \emph{interior boundary}, supported on $\overline{\partial^*\Omega} \cap \mathrm{int}(M)$;
    \item $\partial_F \Omega := \partial \llrr{\Omega} \llcorner \partial M$ is the \emph{free boundary part}, supported on $\overline{\partial^*\Omega} \cap \partial M$.
\end{itemize}
\end{definition}

In Sections~\ref{ch:nested-sweepout} and~\ref{ch:no-escape-to-infinity}, we localise a global sweepout to a good set $U \subset M$ in order to apply the min-max theory on a precompact domain. This localisation requires cutting the Caccioppoli sets $\Omega_t$ along the interior boundary $\partial_I U$, which may introduce extra area near $\partial_I U$. To control this, we use the coarea inequality to find cut-off radii where the extra area is small. Throughout we use the following standard fact:

\begin{lemma}[Coarea inequality]\label{lem:coarea}
Let $\Omega \subset M$ be a bounded open set of finite perimeter and $f\colon M \to \mathbb{R}$ a Lipschitz function. Then for almost every $t \in \mathbb{R}$, the level set $f^{-1}(t) \cap \Omega$ has finite $\mathcal{H}^n$-measure, and
\[
\int_{\mathbb{R}} \mathcal{H}^n(f^{-1}(t) \cap \Omega)\, dt \leq \int_\Omega |\nabla f|\, d\mathcal{H}^{n+1}.
\]
\end{lemma}

The min-max procedure produces limit objects not as currents but as varifolds: the pull-tight argument of Section~\ref{ch:convergence-of-min-max} yields a stationary varifold, and the regularity theory (\cite[Theorem~6.2]{LZ17}; see also Chapter~\ref{ch:alternative-regularity}) recovers a smooth minimal hypersurface from it. We recall the basic definitions; standard references include Pitts~\cite{Pit81}, Simon~\cite{Sim84}, Colding--De~Lellis~\cite{CL03}, and De~Lellis--Tasnady~\cite{DLT13}.

\begin{definition}[Varifolds]\label{def:varifold}
Let $\mathrm{Gr}_n(M)$ denote the Grassmannian bundle of (unoriented) $n$-planes over $M$. An \emph{$n$-varifold} on $M$ is a Radon measure on $\mathrm{Gr}_n(M)$. The space of all $n$-varifolds is denoted $\mathcal{V}_n(M)$, equipped with the weak-$*$ topology.

For $V \in \mathcal{V}_n(M)$, the \emph{mass measure} $\|V\|$ is the pushforward of $V$ under the projection $\mathrm{Gr}_n(M) \to M$. The \emph{total mass} is $\|V\|(M)$, and the \emph{support} is $\operatorname{spt}\|V\|$. For an open set $A \subset M$, the \emph{restriction} $V \llcorner A$ is the varifold defined by $V \llcorner A(E) := V(E \cap \mathrm{Gr}_n(A))$.
\end{definition}

\begin{definition}[First variation]\label{def:first-variation}
The \emph{first variation} of $V \in \mathcal{V}_n(M)$ is the linear functional on smooth vector fields defined by
\[
\delta V(X) := \int_{\mathrm{Gr}_n(M)} \operatorname{div}_S X(x) \, dV(x, S),
\]
where $\operatorname{div}_S X(x) := \sum_{i=1}^n \langle \nabla_{e_i} X, e_i \rangle$ for any orthonormal basis $\{e_i\}$ of $S$.
\end{definition}

\begin{definition}[Stationarity]\label{def:stationary}
A varifold $V \in \mathcal{V}_n(M)$ is:
\begin{enumerate}[label=\textup{(\roman*)}]
    \item \emph{stationary in $\operatorname{int}(M)$} if $\delta V(X) = 0$ for all $X \in \mathfrak{X}_c(\operatorname{int}(M))$;
    \item \emph{stationary with free boundary} if $\delta V(X) = 0$ for all $X \in \mathfrak{X}_{\mathrm{tan}}(M)$.
\end{enumerate}
\end{definition}

The pull-tight procedure (Section~\ref{ch:convergence-of-min-max}) requires a metric on the space of varifolds that metrises weak-$*$ convergence and behaves well under localisation to open subsets. This role is played by the $\mathbf{F}$-metric.

\begin{definition}[$\mathbf{F}$-metric]\label{def:F-metric}
The \emph{$\mathbf{F}$-metric} on $\mathcal{V}_n(M)$ is defined by
\[
\mathbf{F}(V, W) := \sup\biggl\{\int f \, dV - \int f \, dW : f \in C_c(\mathrm{Gr}_n(M)),\ |f| \leq 1,\ \operatorname{Lip}(f) \leq 1\biggr\}.
\]
On sets of uniformly bounded mass, $\mathbf{F}$ metrizes the weak-$*$ topology. For an open set $U \subset M$, the \emph{localized $\mathbf{F}$-metric} is $\mathbf{F}_U(V, W) := \mathbf{F}(V \llcorner U, W \llcorner U)$. See \cite[Section~2]{LZ17}.
\end{definition}

%%%%%%%%%%%%%%%%%%%%%%%%%%%%%%%%%%%%%%%%%%%%%%%%%%%%%%%%%%%%%%%%%%%%%%%%%%%%%%%
\subsection{Sweepouts, Width, and Good Sets}\label{sec:sweepouts}
%%%%%%%%%%%%%%%%%%%%%%%%%%%%%%%%%%%%%%%%%%%%%%%%%%%%%%%%%%%%%%%%%%%%%%%%%%%%%%%

The cut-and-paste operations in our proof (splitting, extension, and nested sweepout construction in Sections~\ref{ch:splitting-and-extension}--\ref{ch:no-escape-to-infinity}) are most naturally described using Caccioppoli sets. We therefore adopt a continuous sweepout framework based on Caccioppoli sets (sets of finite perimeter), following the approach of Chambers--Liokumovich~\cite{CL20} and De~Lellis--Tasnady~\cite{DLT13}. In Proposition~\ref{prop:sweepout-cycle} below, we show that every Caccioppoli sweepout gives rise to a continuous map into the space of mod~$2$ relative cycles; this is the bridge that allows us to invoke the min-max machinery of Li--Zhou~\cite{LZ17} (pull-tight, almost minimizing, regularity) developed in the cycle-valued setting. For applications to non-compact manifolds, we consider sweepouts of bounded open sets.

\begin{definition}[Sweepout of $(U, \partial_F U)$]\label{def:sweepout}
Let $U \subset M$ be a bounded relatively open set of finite perimeter. A \emph{sweepout of $(U, \partial_F U)$} is a family of Caccioppoli sets $\{\Omega_t\}_{t \in [0,1]}$ in $M$ satisfying:
\begin{enumerate}[label=\textup{(sw\arabic*)}]
    \item\label{sw:boundary} \textbf{Boundary condition}: $\Omega_0 \cap U = \emptyset$ and $U \subset \Omega_1$;
    \item\label{sw:continuity} \textbf{Continuity}: the map $t \mapsto \mathbf{1}_{\Omega_t}$ is continuous in $L^1(M)$, i.e.,
    \[
    \mathcal{H}^{n+1}(\Omega_t \triangle \Omega_s) \to 0 \quad \text{as } s \to t;
    \]
    \item\label{sw:finite-mass} \textbf{Finite mass}: $\sup_{t \in [0,1]} \mathrm{Per}(\Omega_t) < \infty$;
    \item\label{sw:no-concentration} \textbf{No mass concentration}: for every $x \in M$,
    \[
    \lim_{r \to 0} \sup_{t \in [0,1]} \mathrm{Per}(\Omega_t, B_r(x)) = 0.
    \]
\end{enumerate}
The \emph{height} of the sweepout is $F(\{\Omega_t\}) := \sup_{t \in [0,1]} \mathrm{Per}(\Omega_t)$.
\end{definition}

\begin{definition}[Width]\label{def:width}
The \emph{width} of $(U, \partial_F U)$ is
\[
W(U) := \inf\bigl\{ F(\{\Omega_t\}) : \{\Omega_t\} \text{ is a sweepout of } (U, \partial_F U) \bigr\}.
\]
\end{definition}

The following proposition provides the bridge between the Caccioppoli sweepout framework and the cycle-valued maps used in Li--Zhou~\cite{LZ17}.

\begin{proposition}[From Caccioppoli sweepouts to cycle-valued maps]\label{prop:sweepout-cycle}
Let $\{\Omega_t\}_{t \in [0,1]}$ be a sweepout of $(U, \partial_F U)$. Then the map
\[
\Phi: [0,1] \to \mathcal{Z}_n(M, \partial M;\mathbb{Z}_2), \quad t \mapsto \partial\llrr{\Omega_t}
\]
is continuous in the flat topology $\mathcal{F}$, and satisfies $\mathbf{M}(\Phi(t)) = \mathrm{Per}(\Omega_t)$ for all $t$ by Proposition~\ref{prop:cacc-cycle}. Moreover, $\Phi$ satisfies the no mass concentration condition
\[
\lim_{r \to 0} \sup_{t \in [0,1]} \|\Phi(t)\|(B_r(x)) = 0 \quad \text{for every } x \in M.
\]
\end{proposition}

\begin{proof}
The mass identity and no concentration condition are immediate from Proposition~\ref{prop:cacc-cycle} and~\ref{sw:no-concentration}. For $\mathcal{F}$-continuity, $\llrr{\Omega_t} - \llrr{\Omega_s}$ is a filling of $\Phi(t) - \Phi(s)$ modulo $\partial M$, so $\mathcal{F}(\Phi(t) - \Phi(s)) \leq \mathcal{H}^{n+1}(\Omega_t \triangle \Omega_s) \to 0$ by~\ref{sw:continuity}.
\end{proof}

\begin{definition}[Good sweepout]\label{def:good-sweepout}
A sweepout $\{\Omega_t\}$ of $(U, \partial_F U)$ is a \emph{good sweepout} if there exists $\varepsilon > 0$ such that
\[
\mathrm{Per}(\Omega_t) \leq 5\,\mathbf{M}(\partial_I U) \quad \text{for all } t \in [0, \varepsilon) \cup (1-\varepsilon, 1].
\]
We write $W_g(U)$ for the infimum of heights over all good sweepouts.
\end{definition}

\begin{definition}[Nested sweepout]\label{def:nested-sweepout}
A sweepout $\{\Omega_t\}$ of $(U, \partial_F U)$ is \emph{nested} if $\Omega_s \subset \Omega_t$ for all $0 \leq s \leq t \leq 1$. Equivalently, $\{\Omega_t\}$ arises from a $\partial$-transverse Morse function $f: M \to \mathbb{R}$ via $\Omega_t = f^{-1}((-\infty, c_t))$ for some increasing function $t \mapsto c_t$. We denote the collection of nested sweepouts of $U$ by $\mathcal{S}_n(U)$, and write $W_n(U) := \inf_{S \in \mathcal{S}_n(U)} F(S)$.
\end{definition}

\begin{definition}[Relative sweepout and relative width]\label{def:relative-sweepout}
A \emph{relative sweepout} of $U$ is a sweepout $\{\Omega_t\}$ of $(U, \partial_F U)$ with the additional property that $\partial^*\Omega_t \subset \overline{U}$ for all $t \in [0,1]$. We denote the collection of relative sweepouts by $\mathcal{S}_\partial(U)$, and define the \emph{relative width} as $W_\partial(U) := \inf_{S \in \mathcal{S}_\partial(U)} F(S)$.
\end{definition}

By Proposition~\ref{prop:sweepout-cycle}, each of the above sweepout classes (good, nested, relative) gives rise to an $\mathcal{F}$-continuous map into $\mathcal{Z}_n(M, \partial M;\mathbb{Z}_2)$ satisfying the no mass concentration condition, so the min-max machinery of Li--Zhou~\cite{LZ17} applies to each corresponding width.

\begin{proposition}[Width inequalities]\label{prop:width-inequalities}
Let $U \subset M$ be a bounded open set with smooth boundary. Then:
\begin{enumerate}[label=\textup{(\roman*)}]
    \item $W_\partial(U) \leq W_n(U) \leq W_\partial(U) + \mathbf{M}(\partial_I U)$;
    \item $W_n(U) = W(U)$, where $W(U)$ is the width over all sweepouts;
    \item If $U$ is a good set (Definition~\ref{def:good-set} below), then $W_g(U) = W(U)$.
\end{enumerate}
\end{proposition}

\begin{proof}
Part~(i) follows by restricting a nested sweepout to $\overline{U}$ (giving the lower bound) and extending a relative sweepout by adding $\partial_I U$ at each time (giving the upper bound). Part~(ii) is proved by approximating an arbitrary sweepout by nested ones using the coarea formula; see \cite[Lemma~4.2]{CL20}. Part~(iii) follows from~(ii) by a cut-off argument using the good set condition.
\end{proof}


\begin{definition}[Good set]\label{def:good-set}
A bounded open subset $U \subset M$ of finite perimeter is a \emph{good set} if
\[
\mathbf{M}(\partial_I U) \leq \frac{1}{10}\, W_\partial(U),
\]
where $\partial_I U = \partial \llbracket U \rrbracket \llcorner \mathrm{int}(M)$ is the interior boundary of $U$.
\end{definition}

\begin{lemma}[Existence of good sets]\label{lem:good-set-existence}
Let $(M^{n+1}, g)$ be a complete Riemannian manifold with non-empty boundary and sublinear volume growth, i.e.\
\[
\liminf_{R \to \infty} \frac{\mathcal{H}^{n+1}(B_R(x) \cap M)}{R} = 0 \qquad \text{for some } x \in M.
\]
Then there exists a good set $U \subset M$ with $0 < W_g(U) < \infty$.
\end{lemma}

\begin{proof}
Fix a point $x \in M$ and a radius $r > 0$ such that the local isoperimetric constant
\[
C_I := \inf\bigl\{\mathbf{M}(\partial_I E) : E \subset B_r(x),\ \mathcal{H}^{n+1}(E) = \tfrac{1}{2}\,\mathcal{H}^{n+1}(B_r(x))\bigr\} > 0.
\]
Let $d_x: M \to [0, \infty)$ denote the distance function from $x$. Applying the coarea inequality (Lemma~\ref{lem:coarea}) with $\Omega = M \cap B_S(x)$ and $f = d_x$ (noting $|\nabla d_x| \leq 1$), we obtain
\[
\int_0^S \mathcal{H}^n(d_x^{-1}(R) \cap \mathrm{int}(M))\, dR \leq \mathcal{H}^{n+1}(M \cap B_S(x))
\]
for any $S > 0$. The sublinear growth assumption therefore allows us to pick a regular value $R > r$ such that
\[
\mathcal{H}^n(d_x^{-1}(R) \cap \mathrm{int}(M)) \leq \frac{C_I}{100}.
\]
Set $U := B_R(x) \cap M$. Since $\partial_I U \subset d_x^{-1}(R) \cap \mathrm{int}(M)$, we have $\mathbf{M}(\partial_I U) \leq C_I/100$. Any sweepout of $U$ must bisect volume in $B_r(x)$ at some time, so $W_\partial(U) \geq C_I$. Since $\mathbf{M}(\partial_I U) \leq \frac{1}{10} W_\partial(U)$, the set $U$ is good. Finiteness $W_g(U) < \infty$ follows from the existence of a sweepout by level sets $\{d_x < R'\}$ for $R' \in [0, R]$.
\end{proof}